% !TeX root = ../thesis.tex
\chapter{Diskussion}
\label{chap:diskussion}

Dieses Kapitel interpretiert die empirischen und technischen Befunde aus Kapitel~\ref{chap:ergebnisse} im Kontext des theoretischen Forschungsstands (Abschnitt~\ref{sec:diskussion-interpretation}), beantwortet systematisch die Forschungsfragen der Arbeit (Abschnitt~\ref{sec:diskussion-forschungsfragen}), analysiert die beobachteten Zielkonflikte zwischen den Verfahrensklassen (Abschnitt~\ref{sec:diskussion-tradeoffs}), leitet praktische Handlungsempfehlungen für den Unternehmenseinsatz ab (Abschnitt~\ref{sec:diskussion-praxis}) und reflektiert die methodischen und datenseitigen Limitationen der Untersuchung (Abschnitt~\ref{sec:diskussion-limitationen}).

\section{Interpretation und Einordnung der Ergebnisse in den Forschungsstand}
\label{sec:diskussion-interpretation}

Die explorativen Entwicklungsbefunde dieser Arbeit sind vereinbar mit zentralen Aussagen der aktuellen Literatur zur automatisierten Dokumentenverarbeitung~\cite{Skalicky2022Benchmarks,Fu2025OCRBenchV2,Kim2022Donut}. Zunächst zeigt sich, dass die Extraktion aus realen Unternehmens-Teilnehmerlisten kein isoliertes Texterkennungsproblem ist, sondern durch das Zusammenspiel aus Bildgeometrie, Tabellenstruktur und semantischer Rollenzuordnung bestimmt wird~\cite{OGorman1997DIA,Jaume2019FUNSD}. Während klassische \ac{ocr} (Tesseract~\cite{Smith2007Tesseract}) auf gedruckten Zeilen solide Worttranskriptionen liefert, scheitert eine starre regelbasierte Nachverarbeitung (\texttt{ocr\_rules}, $0/184$ Felder im Dreidokumentlauf), sobald Formulare abweichende Spaltenüberschriften (\emph{Benutzer}, \emph{Kurzzeichen}, \emph{User}) oder kombinierte Namensspalten verwenden.

Im Vergleich der beiden \ac{ocr}-gestützten Sprachmodell-Pipelines (\texttt{ocr\_llm} vs.\ \texttt{paddle\_llm}) bei identischem Sprachmodell (\texttt{Qwen3-4B-Instruct-2507}~\cite{Yang2025Qwen3}) erweist sich die Qualität der vorgeschalteten Erkennungs- und Layoutstufe als ausschlaggebend: \texttt{paddle\_llm} übertrifft \texttt{ocr\_llm} sowohl auf dem gemischten Einzeldokument ($40/72$ vs.\ $29/72$ exakte Felder) als auch im Dreidokumentvergleich ($94/184$ vs.\ $74/184$ exakte Felder; Feldabdeckung $68{,}48\,\%$ vs.\ $49{,}46\,\%$). Dies deckt sich mit den Befunden von van Strien et al.~\cite{vanStrien2020OCRQuality}, wonach nachgelagerte Sprachmodelle starke Auslassungen oder Verzerrungen einer schwachen \ac{ocr}-Stufe nicht aus dem Text allein heilen können -- insbesondere dann nicht, wenn das System zur Vermeidung von Halluzinationen~\cite{Huang2024Hallucination} zu Recht an eine strikte wörtliche Zeilenbelegprüfung gebunden ist.

Gleichzeitig relativieren die Ergebnisse die Annahme einer voraussetzungslosen Überlegenheit rein \ac{ocr}-freier Vision-Language-Modelle~\cite{Kim2022Donut,Wang2024Qwen2VL}. Zwar kann \texttt{Qwen3-VL-4B-Instruct} (\texttt{vlm}) handschriftliche Namen direkt aus dem Seitenbild lesen, die für Tesseract unzugänglich bleiben; auf mehrseitigen Formularen ohne wiederholte Kopfzeile leidet das direkte \ac{vlm} jedoch unter zwei gravierenden Schwächen: Erstens bricht die autoregressive Generierung auf dichten Folgeseiten mehrfach am Ausgabetoken-Limit von $8\,192$~Tokens mit unvollständigem \ac{json} ab (zwei von sechs Seiten im Dreidokumentlauf, vier von $20$~Seiten im Zehndokumentlauf), was den Beobachtungen von Fu et al.~\cite{Fu2025OCRBenchV2} (OCRBench~v2) und Geng et al.~\cite{Geng2023Grammar} entspricht. Zweitens gibt das Modell auf kopflosen Folgeseiten häufig die erste Teilnehmerzeile fälschlich als Tabellenüberschrift (\texttt{identity\_headers}) aus. Erst die hybride Kopplung mit PaddleOCR-Belegen (\texttt{ocr\_vlm}) und die geometrisch-native Ankerprüfung (\texttt{echo\_corroboration}) überwinden diesen Strukturverlust und steigern die exakten Treffer im Dreidokumentlauf von $67$ (\texttt{vlm}) auf $88$ beziehungsweise $96$ (\texttt{ocr\_vlm}).

\section{Zusätzliche Grenzen der Bewertungsgrundlage}
\label{sec:diskussion-validitaet}

Die Dokumentprüfung in Abschnitt~\ref{sec:ergebnisse-validitaet} bestätigt konkrete Unterschiede zwischen Kursstammdaten und sichtbarem Papierinhalt. Sie begründet jedoch keine automatische Bereinigung aller Referenzfehler. Selektiv lesbare, überwiegend gedruckte Prüfwerte können eine Zusatzmetrik deutlich begünstigen; ungeprüfte Fälle bleiben deshalb im ursprünglichen Nenner. Eine unabhängige Zweitannotation und ein vorab festgelegtes repräsentatives Prüfprotokoll sind für eine belastbare dokumentbasierte Leistungsschätzung erforderlich.

ANLS toleriert Schreibabweichungen, bestätigt aber keine Personalnummer oder Person. Das bedingte Ähnlichkeitsmittel kann durch selektive Auslassungen hoch ausfallen; die ANLS-Feldadaption bestraft solche Auslassungen, erfasst jedoch keine Zusatzpersonen. Beide müssen mit Coverage, Exact Recall/F1 und Teilnehmermetriken gelesen werden. Best-Fit nutzt die Referenz und kann die Zuordnung begünstigen. Die konservativen Gegenproben beschreiben diese Sensitivität, beheben aber keine fehlende unabhängige Identitäts-Ground-Truth. Rollenvertauschungen werden in der Primäransicht weiterhin streng bewertet.

Das mehrdimensionale Vorbild DUE und die unterschiedlichen Aufgaben von DocVQA, OCRBench, FUNSD und PubTabNet sprechen für eine aufgabenspezifische Interpretation~\cite{Borchmann2021DUE,Mathew2021DocVQA,Liu2024OCRBench,Jaume2019FUNSD,Zhong2020PubTabNet}. Ohne annotierte Tabellenbäume, Beziehungskanten oder vollständige Seitentranskripte wären TEDS, Beziehungs-F1 oder OCR-CER/WER nicht belastbar. Die implementierte Textauffindbarkeit und Zeilenkonsistenz sind enger begrenzte Diagnosen. Ungepaarte Zusatzwerte bleiben mögliche Fehlzuordnungen oder unbelegte Extraktionen; sie sind nicht automatisch bewiesene Halluzinationen.

Seiten- und Qualitätsgruppen benötigen unabhängige Referenzlokalisierung. Andernfalls würden gerade nicht erkannte Teilnehmer aus den schwierigen Gruppen verschwinden. Die bisher lokalisierten $34$ Personen erlauben nur eine explorative Teilanalyse. Technische JSON-/Schemaquoten beschreiben letzte gespeicherte Antworten und können Cachewiederverwendung enthalten; sie sind keine vollständige Erstversuchs-Erfolgsrate. Wiederholte Seitenversuche, Modell-Retries und kontrollierte Cold-/Warm-Starts dürfen nicht gleichgesetzt werden. Für Laufzeit- und Ressourcenrangfolgen fehlen weiterhin kontrollierte Wiederholungen unter gleicher Hardware-, Cache- und Startkonfiguration.

\section{Beantwortung der Forschungsfragen}
\label{sec:diskussion-forschungsfragen}

Auf Grundlage der durchgeführten Untersuchungen lassen sich die vier Unterfragen sowie die übergeordnete Forschungsfrage wie folgt beantworten:

\begin{enumerate}
  \item \textbf{Zu FF1 (Extraktionsqualität und Fehlerursachen):} Unter den untersuchten Verfahren erreichen \texttt{paddle\_llm} ($R_{\mathrm{exact}} = 51{,}09\,\%$, $F_{1,\mathrm{exact}} = 60{,}65\,\%$ im Dreidokumentlauf; $55{,}56\,\%$ auf \texttt{140829.pdf}) und das Hybridverfahren \texttt{ocr\_vlm} ($R_{\mathrm{exact}} = 47{,}83\,\%$ primär, $52{,}17\,\%$ nach abgesicherter Folgeseiten- und Swap-Bewertung; $58{,}33\,\%$ auf \texttt{140829.pdf}) die höchste Extraktionsqualität. Vor- und Nachname sowie Personalnummern werden wesentlich zuverlässiger extrahiert als Organisationskürzel (bei \texttt{paddle\_llm} $31/46$ Vorname, $26/46$ Nachname, $25/46$ \ac{id}, aber nur $12/46$ Organisation exakt). Die kausale Fehleranalyse weist nach, dass scheinbare Feldverluste neben echten Handschrift-Lesefehlern vor allem durch (a) seitlich gedrehte Scans bei unzureichendem Standard-\ac{osd}, (b) fehlende Spaltenüberschriften auf Folgeseiten, (c) Eingabe- und Ausgabetoken-Abbrüche, (d) zeilenweise wechselnde Namensreihenfolgen und (e) in externen Kursreferenzen geführte, auf dem Papierformular jedoch physisch fehlende \ac{id}-/Abteilungsspalten verursacht werden.
  \item \textbf{Zu FF2 (Laufzeit- und Ressourceneffizienz):} Durch sequenzielle Prozesssteuerung und \ac{bf16}-Präzision mit \texttt{sdpa\_repeated\_kv} lassen sich sämtliche 4B-Sprach- und Bildmodelle stabil innerhalb des \mbox{16-GiB}-\ac{vram}-Budgets ausführen (Spitzenallokation $\le 10\,\text{GiB}$). Während die beiden Text-\ac{llm}-Pipelines \texttt{ocr\_llm} ($\approx 22{,}1\,\mathrm{s}$/Seite) und \texttt{paddle\_llm} ($21{,}1$--$48{,}5\,\mathrm{s}$ pro frischer Seite, $< 0{,}1\,\mathrm{s}$ bei Cache-Treffer, $9{,}05\,\text{GiB}$ \ac{vram}) moderaten Rechenaufwand erfordern, benötigt das Hybridverfahren \texttt{ocr\_vlm} im Median rund $47{,}9\,\mathrm{s}$ pro Seite. Beim direkten \ac{vlm} dominieren einzelne Tokenlimit-Abbrüche ($> 400\,\mathrm{s}$ pro abgebrochener Seite) die Gesamtlaufzeit, während die zusätzlichen PaddleOCR-Belege im Hybridpfad \texttt{ocr\_vlm} auf denselben sechs Seiten solche Endlos-Generierungen verhinderten.
  \item \textbf{Zu FF3 (Wirtschaftlichkeit und Nachbearbeitungsaufwand):} Da auf realen, gemischten Listen selbst die besten Verfahren derzeit einen strikten Teilnehmer-Recall von knapp $48\,\%$ und eine Feldabdeckung von knapp $69\,\%$ erreichen, ist ein vollautomatischer Dunkelverbuchungsbetrieb ohne menschliche Nachkontrolle fachlich ausgeschlossen. Der wirtschaftliche Mehrwert entsteht im \emph{Human-in-the-Loop}-Betrieb: Die Pipeline richtet gedrehte Seiten automatisch aus, strukturiert bis zu zwei Drittel der Felder vor, hebt unsichere Werte (\texttt{state="uncertain"}) sowie den fünfstufigen Feldstatus im Dashboard hervor und verlinkt jedes Feld direkt mit dem Original-\ac{pdf}-Ausschnitt, wodurch sich die manuelle Erfassungszeit $t_{\text{manuell}}$ auf eine deutlich kürzere Prüf- und Ergänzungszeit $t_{\text{Prüfung}}$ reduziert.
  \item \textbf{Zu FF4 (Praktische Einsatzfähigkeit und Robustheit):} Der praktische Einsatz erfordert zwingend die im Framework implementierten Schutzschichten: die zweistufige Seitenorientierung (\ac{osd} + Vierlagen-\ac{ocr}), die zerstörungsfreie Hintergrundnormalisierung für graue Tabellenköpfe, die geometrisch geprüfte Spaltenvererbung ($\Delta x_{\max} \le 0{,}035$), die zeilenindividuelle Namensaufteilung und die verlustfreie Input-Kompaktierung. Durch die strikte Offline-Kapselung und das isolierte Transferpaket \texttt{software/api\_firmen\_pc/} sind sowohl der lokale \ac{gpu}-Betrieb als auch die spätere Anbindung an einen unternehmensinternen \ac{api}-Endpunkt datenschutzkonform realisierbar.
\end{enumerate}

\section{Zielkonflikte (Trade-offs) zwischen den untersuchten Ansätzen}
\label{sec:diskussion-tradeoffs}

Die experimentellen Ablationen legen vier grundlegende Zielkonflikte offen, die bei der Konfiguration eines produktiven Extraktionssystems abgewogen werden müssen:
\begin{itemize}
  \item \textbf{Recall vs.\ Precision bei Handschrift und Zellergänzung:} Sowohl die Beibehaltung unvollständiger Handschriftzeilen (\texttt{preserve\_partial\_fields: true}: exakte Felder $29 \rightarrow 38$, aber falsche Felder $7 \rightarrow 15$) als auch die geometrische Ergänzung fehlender \ac{id}-/Organisationszellen (\texttt{complete\_observed\_cells}: exakte Felder $38 \rightarrow 40$, falsche Felder $15 \rightarrow 20$) und die Input-Kompaktierung ($122 \rightarrow 133$ exakt, $36 \rightarrow 44$ falsch) steigern zwar den Recall und die Feldabdeckung, senken jedoch die exakte Precision, da in schwierig lesbaren Bereichen zusätzliche unvollständig oder fehlerhaft transkribierte \ac{ocr}-Zellen in die Ausgabe gelangen. Die Kennzeichnung dieser Ergänzungen als \texttt{state="uncertain"} löst diesen Zielkonflikt für die menschliche Nachprüfung auf, indem sie dem Prüfpersonal gezielt die unsicheren Kandidaten anzeigt.
  \item \textbf{Bedingte Zeichenähnlichkeit ($S_{\mathrm{detected}}$) vs.\ Vollständigkeit ($\mathrm{Coverage}$):} Ein konservatives Verfahren wie \texttt{ocr\_llm}, das auf schwierigen Handschriftseiten häufig leere Antworten liefert, erzielt auf den wenigen einfachen Seiten eine sehr hohe bedingte Levenshtein-Ähnlichkeit ($93{,}75\,\%$), lässt aber über die Hälfte aller Referenzfelder unentdeckt. Produktiv ist eine höhere Abdeckung (\texttt{paddle\_llm}: $68{,}48\,\%$) mit ausgewiesenen Prüfhinweisen meist wertvoller als ein selektiv hoher Durchschnitt auf wenigen Feldern.
  \item \textbf{Belegstrenge Text-\ac{llm}-Pipeline vs.\ visuelle Hybridflexibilität:} \texttt{paddle\_llm} (PaddleOCR + Text-\ac{llm}, unabhängig vom \ac{vlm}) beschränkt die Übernahme auf protokollierte Literal- und Positionsbelege und ist bei schwer lesbarer Handschrift durch die Zeichenerkennung von \texttt{PP-OCRv5} begrenzt. Diese Belege bestätigen keine Korrektheit gegenüber dem Original. Umgekehrt kann \texttt{ocr\_vlm} (Kombination aus PaddleOCR und \ac{vlm}) visuelle Strukturen und Handschrift direkt im Bild deuten, erfordert aber harte nachgelagerte Filter (wie die Sperre bei Header-Zell-Längenkonflikten und die native Ankerprüfung \texttt{echo\_corroboration}), um Fehlzuordnungen von Unterschriften in Organisationsspalten zu verhindern.
\end{itemize}

\section{Praktische Einordnung und Empfehlungen für den Einsatz bei BMW}
\label{sec:diskussion-praxis}

Für die Überführung des Prototyps in den laufenden Arbeitsablauf bei der BMW Group, um künftig neu eingehende Teilnehmerlisten automatisiert in das Schulungssystem zu übertragen, empfiehlt sich ein gestuftes Vorgehen:
\begin{enumerate}
  \item \textbf{Vorab-Aufbereitung und automatische Ausrichtung neuer Scans:} Zunächst sollten neu eingescannte Teilnehmer-\acp{pdf} über \texttt{vorbereiten.py} anhand der BMW-Dateinamenskonvention nach ihrer sechsstelligen Schulungskennung geordnet und über die zweistufige automatische Orientierung (\texttt{pdf.prepare}) ausgerichtet werden. Seiten, bei denen der Vierlagen-\ac{ocr}-Vergleich wegen extrem geringer Textmenge keinen eindeutigen Sieger liefert (\texttt{orientation\_status="review\_required"}), sind vorab in der Kontaktansicht manuell zu sichten.
  \item \textbf{Wahl der Extraktionsarchitektur:} Steht die lokale GPU-Arbeitsstation mit \mbox{16\,GiB} \ac{vram} zur Verfügung, ist die Kombination aus \textbf{\texttt{paddle\_llm}} (als schnellere, belegstrenge Text-\ac{llm}-Pipeline) und \textbf{\texttt{ocr\_vlm}} (als Hybrid aus PaddleOCR und \ac{vlm} für Seiten mit hohem Handschriftanteil) eine anhand der bisherigen Entwicklungsbefunde zu pr\"ufende Konfiguration; eine allgemeine Leistungs\"uberlegenheit ist noch nicht nachgewiesen. Soll die Verarbeitung hingegen direkt auf einem regulären BMW-Firmen-PC über einen internen \ac{vlm}-Endpunkt (\texttt{software/api\_firmen\_pc/}) erfolgen, muss dort entweder vorab das ausgerichtete Seitenmanifest übergeben oder lokal Tesseract installiert und \texttt{pdf.auto\_orient: true} im \ac{api}-Profil aktiviert werden.
  \item \textbf{Gestützte Nachbearbeitung im Dashboard:} Für die Sachbearbeitung bieten das Offline-Dashboard (insbesondere der Einzelbewertungs- und Filterbereich~G sowie die Fehleransicht \texttt{extraction\_errors.html}) und der in Eingabereihenfolge sortierte Excel-Export (\texttt{participants.xlsx}) eine unmittelbare Arbeitsgrundlage: Als \texttt{read} und hochkonfident eingestufte Zeilen können schnell abgeglichen werden, während sich die manuelle Korrektur auf die als \texttt{uncertain}, \emph{Abweichend} oder \emph{Zuordnung unklar} markierten Felder konzentriert.
\end{enumerate}

\section{Limitationen der Untersuchung}
\label{sec:diskussion-limitationen}

Bei der Interpretation der Arbeit sind fünf methodische und empirische Limitationen zu beachten:
\begin{itemize}
  \item \textbf{Spezifität des Anwendungsfalls bei BMW:} Sämtliche untersuchten Formularvorlagen, Spaltenkonventionen (\emph{Benutzer}, \emph{Benutzername}, \emph{Kurzzeichen}), Dateinamensstrukturen und Referenzexporte stammen aus dem internen Schulungswesen der BMW Group. Da keine Dokumentbestände oder Prozessdaten anderer Unternehmen erhoben wurden, erhebt die Arbeit keinen Anspruch auf eine pauschale Übertragbarkeit der numerischen Ergebnisse auf fremde Unternehmensorganisationen.
  \item \textbf{Entwicklungsstatus der bisher ausgewerteten Stichproben:} Die detailliert untersuchten realen Dokumente (ein zweiseitiges Dokument mit $72$~Feldern, drei Dokumente mit $184$~Feldern und zehn Dokumente mit $620$~Feldern) wurden im Zuge der Fehleranalyse und Pipeline-Reparatur mehrfach gesichtet. Sie belegen die technische Wirksamkeit der implementierten Reparaturen, ersetzen jedoch noch keine unabhängige konfirmatorische Evaluation auf unberührten Pilot- und Testgruppen~\cite{Cawley2010SelectionBias}. Aufgrund der geringen Anzahl unabhängiger Gruppen ($G < 10$) in diesen Entwicklungsläufen sind noch keine belastbaren Cluster-Bootstrap-Konfidenzintervalle ausgewiesen~\cite{Field2007Cluster}.
  \item \textbf{Eigenschaften der externen Kursreferenz:} Da die historischen Excel-Kurslisten weder Seitenkoordinaten noch eine Garantie dafür bieten, dass alle Stammdatenfelder (insbesondere \emph{User~\ac{id}} und \emph{Organization}) auf dem jeweiligen Papierformular tatsächlich als Spalte existierten oder ausgefüllt wurden, kann der Referenznenner $R$ auch dokumentbedingt nicht extrahierbare Informationen enthalten. Die Zusatzbewertung auf gepr\"uften lesbaren Originalfeldern ersetzt ohne repr\"asentative Vollannotation keinen Gesamtscore. Ebenso ist das Best-Fit-Matching eine referenzgestützte Paarung und keine unabhängige Identitätsverifikation.
  \item \textbf{Fehlende lokalisierte Handschrift-Transkripte und Signatur-Ground-Truth im Kursvergleich:} Ohne manuell annotierte Wort- und Zellen-Bounding-Boxes lässt sich auf nicht erkannten Regionen nicht für jedes einzelne Feld beweisen, ob die Textdetektion oder die Zeichentranskription scheiterte. Zudem enthalten die Excel-Kurslisten keine Spalte für das Vorhandensein einer Unterschrift, sodass die quantitative Evaluation von \texttt{signature\_present} auf das separate manuelle Seitenreferenz-Protokoll (\texttt{ocr-reference}) angewiesen ist.
  \item \textbf{Kein kontrollierter Kaltstart-Laufzeitbenchmark:} Da mehrere Entwicklungsläufe bestehende \ac{ocr}- oder Strukturcaches nutzten oder parallel zu Testläufen stattfanden, beschreiben die berichteten Zeiten das reale Laufverhalten des jeweiligen Laufs, stellen jedoch noch keine unter identischen Kaltstartbedingungen wiederholte Latenzmessreihe dar.
\end{itemize}
